% Part: lambda-calculus
% Chapter: lambda-definability
% Section: truth-values

\documentclass[../../../include/open-logic-section]{subfiles}

\begin{document}

\olfileid{lam}{ldf}{tvr}
\olsection{সত্যমান ও সম্বন্ধ}

বিশুদ্ধ ল্যাম্বডা ক্যালকুলাসে সত্যমানগুলিকে নিচের মতো সাংকেতিত করা যায়:
\begin{align*}
  \fn{true} & \ident \lambd[x][\lambd[y][x]]\\
  \fn{false} & \ident \lambd[x][\lambd[y][y]]
\end{align*}

সত্যমানগুলিকে \emph{নির্বাচক} হিসেবে উপস্থাপন করা হয়; অর্থাৎ, এমন
অপেক্ষক হিসেবে, যা দুটি আর্গুমেন্ট গ্রহণ করে এবং তাদের একটি ফেরত দেয়।
সত্যমান $\fn{true}$ তার প্রথম আর্গুমেন্টটি এবং $\fn{false}$ তার দ্বিতীয়
আর্গুমেন্টটি নির্বাচন করে। যেমন, $\fn{true}\, M N$ সর্বদা $M$-এ হ্রাস
পায়, আর $\fn{false}\, M N$ সর্বদা~$N$-এ হ্রাস পায়।

\begin{defn}
কোনো সম্বন্ধ $R \subseteq \Nat^k$-কে !!{lambda definable} বলি, যদি এমন
একটি পদ~$R$ থাকে যাতে
\begin{align*}
  R\, \num{n_1} \dots \num{n_k} & \bred \fn{true}
  \intertext{যখন $R(n_1, \dots, n_k)$ সত্য, এবং}
  R\, \num{n_1} \dots \num{n_k} & \bred \fn{false}
\end{align*}
অন্যথায়।
% BN-SRC-308 TARGET-CORRECTION-BEGIN
\par\smallskip\noindent\textnormal{\small\textbf{উৎস-সংশোধন (BN-SRC-308):} হিমায়িত সংজ্ঞা সম্বন্ধটিকে প্রথমে $R\subseteq\Nat^n$ বললেও উভয় হ্রাস এবং সত্যতার শর্তে আর্গুমেন্টগুলি $n_1,\dots,n_k$। বাংলা পাঠে ঘোষিত স্থানসংখ্যাকে ওই সর্বত্র-ব্যবহৃত~$k$-এর সঙ্গে মিলিয়ে $R\subseteq\Nat^k$ করা হয়েছে।} % BN-SRC-308 TARGET-CORRECTION-END
\end{defn}

যেমন, $\fn{IsZero} = \{0\}$ সম্বন্ধটি $0$-এর ক্ষেত্রে, এবং কেবল
$0$-এর ক্ষেত্রেই সত্য; এটি নিচের পদটির দ্বারা !!{lambda definable}:
\[
  \fn{IsZero} \ident \lambd[n][n (\lambd[x][\fn{false}])\, \fn{true}].
\]
এটি কীভাবে কাজ করে? চার্চ সংখ্যাপদগুলি পুনরাবর্তক হিসেবে সংজ্ঞায়িত
(এমন অপেক্ষক, যা তার প্রথম আর্গুমেন্টকে দ্বিতীয়টির উপর $n$ বার প্রয়োগ
করে), তাই প্রারম্ভিক মান $\fn{true}$ নিই এবং পুনরাবৃত্তির প্রতিটি ধাপে
আগের ধাপের ফল যা-ই হোক না কেন, $\fn{false}$ ফেরত দিই। এই ধাপটি
প্রারম্ভিক মানটির উপর $n$ বার প্রয়োগ হবে; ফলে ধাপটি একবারও প্রয়োগ না
হলে, অর্থাৎ $n = 0$ হলেই কেবল ফল হবে $\fn{true}$।

সত্যমানের এই উপস্থাপনের ভিত্তিতে আরও কয়েকটি সত্যমান-অপেক্ষক সংজ্ঞায়িত
করা যায়। এখানে দুটি দেওয়া হলো---নঞর্থকরণ ও সংযোজনের উপস্থাপন:
\begin{align*}
  \fn{Not} & \ident \lambd[x][x\, \fn{false}\, \fn{true}]\\
  \fn{And} & \ident \lambd[x][\lambd[y][xy \,\fn{false}]]
\end{align*}
``$\fn{Not}$'' অপেক্ষকটি একটি আর্গুমেন্ট গ্রহণ করে এবং আর্গুমেন্টটি
$\fn{false}$ হলে $\fn{true}$, আর আর্গুমেন্টটি~$\fn{true}$ হলে
$\fn{false}$ ফেরত দেয়। ``$\fn{And}$'' অপেক্ষকটি আর্গুমেন্ট হিসেবে দুটি
সত্যমান গ্রহণ করে এবং উভয় আর্গুমেন্ট~$\fn{true}$ হলে এবং কেবল তখনই
$\fn{true}$ ফেরত দেওয়ার কথা। সত্যমানগুলিকে নির্বাচকরূপে উপস্থাপন করা
হয় (উপরে যার বর্ণনা আছে), তাই $x$ সত্যমান হলে এবং তাকে দুটি
আর্গুমেন্টের উপর প্রয়োগ করলে, $x$ যদি $\fn{true}$ হয় তবে ফল হবে প্রথম
আর্গুমেন্ট, আর অন্যথায় দ্বিতীয়টি। এবার $\fn{And}$ তার দুই আর্গুমেন্ট
$x$ ও~$y$ নিয়ে প্রথম আর্গুমেন্ট~$x$-এর কাছে $y$ ও~$\fn{false}$ পাঠায়।
$x$-কে সত্যমান ধরে নিলে, $x$ যদি $\fn{true}$ হয় তবে ফল~$y$-তে, আর
$x$ যদি $\fn{false}$ হয় তবে $\fn{false}$-এ মূল্যায়িত হবে; ঠিক এটিই
চাই।

লক্ষ করো, এখানে আমরা ধরে নিচ্ছি যে $\fn{And}$-এর আর্গুমেন্ট হিসেবে কেবল
সত্যমানই ব্যবহৃত হয়। অন্য পদ দিলে ফলটি (অর্থাৎ, স্বাভাবিক রূপের অস্তিত্ব
থাকলে সেই স্বাভাবিক রূপ) সত্যমান না-ও হতে পারে।

\begin{prob}
  সত্যমানগুলিকে $\lambd$-পদ হিসেবে সাংকেতিত করে অন্তর্ভুক্তিমূলক ও
  বর্জনমূলক বিয়োজনের সত্যমান-অপেক্ষক উপস্থাপনকারী $\fn{Or}$ ও~$\fn{Xor}$
  অপেক্ষকদুটি সংজ্ঞায়িত করো।
\end{prob}

\end{document}
